% 摘要后的正文首页按示例版式上移一级标题，使标题落在约 25--30 mm 的页顶区域。
\vspace*{-20pt}
\section{问题重述}

\subsection{问题背景}

在自主移动机器人、无人车等智能装备的导航与定位任务中，单一传感器定位方式易受环境干扰、自身硬件限制等因素影响，存在定位精度不足、鲁棒性差、数据频率单一等问题，难以满足复杂场景下机器人实时高精度定位的工程需求。实际应用中，通常采用多源异构传感器组合定位方案，通过融合不同原理、不同频率的定位数据，实现定位精度、稳定性与实时性的协同提升。

\subsection{问题提出}

现有机器人实际定位场景中的两类核心定位数据，分别为定位方式 1（输出频率 4 Hz，采样周期 0.25 s）与定位方式 2（输出频率 5 Hz，采样周期 0.2 s）。两类定位数据由不同采集设备获取，存在起始时间不同步、采样频率不同、数据含随机噪声问题，直接融合会导致位置信息错位、轨迹失真，无法输出高频率、高精度的连续定位结果。

请根据上述背景和附件数据，完成以下建模任务。射击与拍照约束在问题四后直接列出，便于后文统一引用。

{\heiti 问题一：}附件 1 两类数据在采集过程中无噪声影响，但存在设备开机的先后不同，请建立两类数据的时间对齐模型，并给出两种方式的时间偏差及 10 Hz 的位置轨迹。

{\heiti 问题二：}附件 2 两类数据存在随机测量噪声和固定的系统偏差，请建立数据对齐与融合模型，并给出两种方式的时间偏差和系统偏差，以及 10 Hz 的位置轨迹。

{\heiti 问题三：}附件 3 两类数据是实际测量数据，请判断是否存在系统偏差，在此基础上完成问题二的任务。

{\heiti 问题四：}机器人按照附件 3 中的轨迹运动，在运动过程中需要对沿途目标点执行“模拟射击”或“拍照扫描”任务，目标点坐标见“\texttt{附件 4.xlsx}”。请对任务目标进行优化设计，在下述约束下尽可能多地完成射击和拍照任务，并将结果填入“\texttt{result.xlsx}”中（不改变附件中的红色文字内容）。

{\heiti （1）射击任务约束：}单次射击命中率为 85\%；射击时刻机器人与目标间的距离 $d$ 应满足 $d\in[5\,\mathrm{m},30\,\mathrm{m}]$；射击时刻机器人的线速度 $v$ 应满足 $v\leq 2\,\mathrm{m/s}$；射击时刻机器人的加速度 $a$ 应满足 $a\leq 1.5\,\mathrm{m/s^2}$；机器人射击前需要 1.5 s 的时间进行瞄准校准，在这段时间内，上述约束条件均应满足。

{\heiti （2）拍照任务约束：}对每个拍照目标，要求尽量多的从不同角度进行拍摄；每次拍照距离 $d$ 需满足 $d\in[10\,\mathrm{m},40\,\mathrm{m}]$；同一目标的不同拍照角度之间，方向角差异至少 $60^\circ$；为了保持拍照时的相机稳定，机器人的线速度 $v$ 应满足 $v\leq 1.5\,\mathrm{m/s}$，加速度 $a$ 应满足 $a\leq 1.5\,\mathrm{m/s^2}$；机器人拍照前需要 0.5 s 的时间进行相机对准，这段时间内，上述约束条件均应满足。

\section{问题分析}

\subsection{问题一分析}
问题一中，两种定位方式均为无噪声采样，因此可将观测数据视为来自同一真实连续轨迹的两组离散采样点集。两组数据由于设备开机时刻不同，其时间戳之间存在固定偏移，可分别采用三次样条构造连续位置函数并提取速度特征，利用归一化互相关在全局范围内粗估时间偏差，再以公共时间区间内的位置均方误差为目标，采用 Brent 法进行一维精细搜索。完成时间对齐后，将两组采样映射至统一时间轴并融合重建，得到 10 Hz 位置轨迹。

\subsection{问题二分析}
问题二的两组定位观测同时包含随机测量噪声、固定时间偏差和固定空间偏差。时间错位会改变跨设备观测的对应关系，空间偏差又会改变位置残差的中心，两类偏差在配准目标中相互耦合。为保证不同时间候选使用相同样本并降低局部异常残差的影响，可在固定公共网格上构造 Huber 剖面目标，将二维空间偏差作为给定时间偏差下的伴随估计量。完成时空标定后，根据两种定位方式的测量噪声确定观测协方差，将原始 4 Hz、5 Hz 事件按校正时刻异步更新，再利用 RTS 平滑重建 10 Hz 状态轨迹。

\subsection{问题三分析}
问题三的两组实际定位观测除时间错位和随机测量误差外，是否存在需要校正的系统偏差尚不明确。位置残差具有时序相关性，普通样本均值及独立样本检验不足以支撑偏差判定，需在稳健时间配准后估计残差均值的长期协方差，并联合统计显著性、工程效应、移动块 Bootstrap 置信区间和偏差趋势确定状态维度。根据判定结果选择无偏差、常量偏差或随机游走偏差状态，再完成异步融合；所得 10 Hz 位置、速度、加速度及模型内不确定性共同作为问题四任务约束计算的运动输入。

\subsection{问题四分析}
问题四是在问题三固定轨迹上安排射击和拍照任务。每个候选任务是否可行，不仅取决于执行时刻的距离、速度和加速度，还取决于执行前完整准备区间内上述条件是否始终成立；同一拍照目标的多次任务还存在圆周角差约束。可先识别连续可行窗口并生成代表候选，再按最早完成准则给出任务数的可行下界，最后建立含资源互斥、射击目标唯一性和拍照角差冲突的0--1规划，以任务总数为第一目标、约束裕度为次级目标完成精确求解，不引入题目未规定的固定任务总数。
